\documentclass[12pt,a4paper]{article}
\usepackage{xeCJK}                    % 中文支持（XeLaTeX专用）
\setCJKmainfont{SimSun}               % 使用系统宋体
\usepackage{amsmath,amssymb,amsthm}    % 数学公式
\usepackage{graphicx}                  % 插入图片
\usepackage{booktabs}                   % 表格
\usepackage{algorithm}                  % 算法
\usepackage{algorithmic}                % 算法伪代码
\usepackage{listings}                   % 代码
\usepackage{xcolor}                     % 颜色
\usepackage{hyperref}                   % 超链接
\usepackage{float}                      % 浮动体控制

% 超链接设置
\hypersetup{
    colorlinks=true,
    linkcolor=blue,
    citecolor=green
}

\title{LaTeX高级功能测试（XeLaTeX版本）}
\author{测试}
\date{\today}

\begin{document}

\maketitle

\tableofcontents  % 目录
\newpage

\section{数学公式测试}

\subsection{行内公式}

爱因斯坦质能方程：$E = mc^2$，其中$E$是能量，$m$是质量，$c$是光速。

\subsection{独立公式}

高斯积分：
\begin{equation}
\int_{-\infty}^{\infty} e^{-x^2} dx = \sqrt{\pi}
\label{eq:gaussian}
\end{equation}

矩阵：
\begin{equation}
\mathbf{A} = \begin{pmatrix}
a_{11} & a_{12} \\
a_{21} & a_{22}
\end{pmatrix}
\label{eq:matrix}
\end{equation}

\subsection{多行公式}

\begin{align}
f(x) &= x^2 + 2x + 1 \\
     &= (x+1)^2 \\
     &= (x+1)(x+1)
\label{eq:align}
\end{align}

\section{表格测试}

\begin{table}[H]
\centering
\caption{测试表格}
\label{tab:test}
\begin{tabular}{lcc}
\toprule
项目 & 数值1 & 数值2 \\
\midrule
项目A & 100 & 200 \\
项目B & 150 & 250 \\
项目C & 200 & 300 \\
\bottomrule
\end{tabular}
\end{table}

\section{中文测试}

这是一个包含中文的段落。LaTeX可以很好地处理中文内容，包括各种标点符号：，。！？；：""''（）【】。

\section{列表测试}

\subsection{无序列表}

\begin{itemize}
    \item 第一项
    \item 第二项
    \begin{itemize}
        \item 子项1
        \item 子项2
    \end{itemize}
    \item 第三项
\end{itemize}

\subsection{有序列表}

\begin{enumerate}
    \item 第一步
    \item 第二步
    \begin{enumerate}
        \item 子步骤1
        \item 子步骤2
    \end{enumerate}
    \item 第三步
\end{enumerate}

\section{引用测试}

引用公式\eqref{eq:gaussian}和表格\ref{tab:test}。

\section{测试完成}

如果这个PDF能正常生成，说明LaTeX环境配置成功！✅

\end{document}

